\subsection{Publication Architectures}
\label{sec:protocols}

We compare five publication architectures that share the same 64-byte
record layout and consumer logic. They differ in two design choices: the
transport from producer memory to monitor memory, and the primitive that
binds the field updates into one logical publication.

\begin{table}[t]
\centering
\scriptsize
\setlength{\tabcolsep}{4pt}
\begin{tabular}{@{}lll@{}}
\toprule
Architecture       & Transport         & Primitive \\
\midrule
\unsync            & Coherent line      & None \\
\seqlock           & Coherent line      & Odd-even version \\
\dblbuf            & Coherent slot pair & Generation counter \\
\dmanaive          & DMA-pulled mirror  & None \\
\dmaseqlock        & DMA-pulled mirror  & Odd-even version on mirror \\
\bottomrule
\end{tabular}
\caption{Five publication architectures. The transport moves the record;
the primitive determines whether the record is published as a consistent
snapshot.}
\label{tab:protocols}
\end{table}

The first three designs use direct coherent sharing. \unsync writes and
reads the record with no publication discipline. \seqlock adds an
odd-even version counter so the reader can detect an update in progress
and retry. \dblbuf writes into one of two slots and then publishes a
monotonic generation counter; the reader retries only when the writer
laps it twice during one copy. This generation counter is important. A
binary published-index flag cannot detect the double-lap case because the
flag returns to its original value.

The last two designs use explicit transfer. A custom gem5 DMA pull
engine periodically copies one 64-byte record from producer memory into a
monitor-local mirror page. \dmanaive reads that mirror with no handshake,
so it can tear both when the DMA engine samples the producer mid-write
and when the monitor samples the mirror mid-copy. \dmaseqlock uses the
same transport but validates the mirror with a sequence-lock check. In
the simulator this composition assumes that one DMA pull is observed as
one logical 64-byte mirror update. We treat that as an explicit modeling
assumption rather than as a claim about all hardware DMA engines.
